\section{The locked triple and its common symmetry}
\label{sec:locked}

All subspace dimensions in this paper are vector dimensions.  Thus a
block \(B\in\cB_i\) is a four-dimensional subspace of \(\FF_2^8\) and
contains \(15\) projective points.  The free block and point lattices are
\[
  M_i^{\ZZ}=\ZZ[\cB_i],
  \qquad
  P^{\ZZ}=\ZZ[\cP].
\]
We use superscripts \(\QQ\) and \(k\) for scalar extension.  The integral
and rational relation modules are
\[
  K_i^{\ZZ}=\ker(A_i:M_i^{\ZZ}\to P^{\ZZ}),
  \qquad
  K_i=K_i^{\ZZ}\otimes\QQ.
\]
Since \(A_i\) has full row rank,
\[
 \dim_{\QQ}K_i=66\,929-255=66\,674.
\]

The common group is generated by elements \(a,b,c\) satisfying
\[
 a^{17}=b^4=c^3=1,\qquad bab^{-1}=a^4,
 \qquad [a,c]=[b,c]=1.
\]
Set \(F=\langle a,b\rangle\cong C_{17}\rtimes C_4\).  The action of
\(\GammaG=F\times\langle c\rangle\) preserves every block collection and
the point set, so every incidence map \(A_i\) is
\(\GammaG\)-equivariant.

\subsection{The eight rational types}

The rational irreducibles can be constructed without choosing a matrix
model.  The quotient \(F/C_{17}\cong C_4\) has the three rational types
\(\mathbf1,\eps,\iot\), of dimensions \(1,1,2\).  Here \(\eps\) is the
order-two character and \(\iot\) is the rational form of the pair of
order-four characters.  The group \(C_3\) has rational types
\(\mathbf1,\om\), where \(\om\) is the two-dimensional rational form of
the two nontrivial complex characters.

For a nontrivial complex character \(\chi\) of \(C_{17}\), the \(C_4\)
action has an orbit of size four, and
\(\Ind_{C_{17}}^F\chi\) has complex degree four.  The four such induced
characters form one Galois orbit.  Their rational form is denoted
\(\rh\) and has dimension \(16\).  Tensoring with the nontrivial
\(C_3\)-orbit gives \(\rh\om\), of dimension \(32\).

Thus the rational types and dimensions are
\[
\begin{array}{c|rrrrrrrr}
V&\mathbf1&\eps&\iot&\om&\eps\om&\iot\om&\rh&\rh\om\\
\midrule
\dim_{\QQ}V&1&1&2&2&2&4&16&32.
\end{array}
\]
For completeness, over \(\CC\) the first six rational types account for
all \(12\) linear characters of \(\GammaG\).  The type \(\rh\) accounts
for four degree-four characters and \(\rh\om\) for the remaining eight.
Consequently
\[
 12\cdot1^2+12\cdot4^2=204,
\]
so the list exhausts the complex irreducibles and hence the rational
types.

\begin{proposition}[Point module]
\label{prop:point-module}
The rational point permutation module is
\[
\QQ[\cP]\cong
3\mathbf1\oplus2\eps\oplus3\om\oplus2\eps\om
\oplus5\rh\oplus5\rh\om.
\]
In particular, neither \(\iot\) nor \(\iot\om\) occurs.
\end{proposition}

\begin{proof}
For each of the \(24\) conjugacy classes, the permutation character equals
the number of fixed projective points.  Exact inner products with the
Galois sums constructed above give the displayed multiplicities.  Their
dimension-weighted sum is
\[
3+2+6+4+80+160=255,
\]
which checks the full point dimension.  The class values and inner
products are included in the electronic representation certificate
described in \cref{app:certificates}.
\end{proof}

\begin{remark}
Only the locked finite triple is considered here.  No family
classification is inferred from its common group or its character table.
\end{remark}
